\section{Week in Review}
\sectionkicker{WEEKLY SYNTHESIS}
\begin{wideflow}
{\Large\bfseries 今週は、公開の条件まで動いた}\par
\smallskip
{\color{SurveyMuted}構造効率の公開群と読者の関心、道具面の組み替え、旗艦の看板と道具内配備、動かす土台と地域公開、実験室の証拠と二つの統治が同時に前景化した。}\par
\end{wideflow}

\subsection*{今週何が変わったか}
\begin{enumerate}[leftmargin=1.6em]
\item \textbf{公開の重みが構造で並んだ。} GLM-5.3-Flash、Qwen3.8-Flash-Next、Hy4 preview が同じ週に現れたが、動作規模も利用条件も preview の位置づけもモデルごとに確かめる必要がある。Granite 4.2の道具型訓練の広がりも符号寄りの前提とともに読む。
\item \textbf{道具の面が組み替わった。} 会話の続きの持ち運び、編集画面への置き場の取り込み、会話の文脈化、置き場と審査、OS層の見張りが並び、便利さと預け先の問いが同時に前に出た。
\item \textbf{看板と土台と統治が分かれた。} 旗艦の測定値は出どころと留保ごと受け止め、動かす土台は出荷記録で固め、実験室の証拠と二つの統治の形は対照として並べた。
\end{enumerate}

\subsection*{なぜそれらが一緒に重要なのか}
公開の構造効率は、対応する道具面と動かす土台があって初めて運用へ届く。旗艦の看板は測定の出どころを切り離さない読みで確かめ、安全の証拠と統治の形は等価にせず境界ごと並べる必要がある。今週の変化は、重みを手に入れる入口、動かし方、確かめ方、預け方の条件が別々の話ではなく、利用条件を構成する連鎖として見えるようになった点にある。

この連鎖は製品間の直接的な互換性を意味しない。各公開、報道、一次記録は固有の主体・環境・測定条件を持つ。動作比率の小ささを共通効率へ、利用条件の表示を共通許諾へ、速度や点数を横断的な順位へ自動的に変換してはいけない。

\begin{communitynote}[Weekly Community Movement --- context only]
X上の話題では、公開の重み群への関心と道具面への関心が並んで見られた。これは今週の関心の広がりを示す背景であり、仕様や性能、提供条件、互換性を証明するものではない\cite{w35community}。
\end{communitynote}

\subsection*{次に何を見るべきか}
次週以降は、この連鎖のどこがさらに具体化するかを追う。GLMの重み置き場、Qwenの利用条件の条文、Hy4の一次記録、AHPやOriginやCopilotの一次文書、Harnessの導入実績、grithの適用範囲、Fable 5やGPT-5.6の一次測定、vLLMの適用、OpenThaiの再現、自動整合の追試、二つの統治の手引きはまだ詳細化されていない。さらに、実験室や提供側の数値が、それぞれの条件から外へ一般化できるかを確認する必要がある。

\begin{claimboundary}[この総括の読み方]
各項目の公開条件、報告主体、測定環境は同じではない。提供元が示す速度、報道が伝える仕様、実験室が報告する評価、話題で見られた関心は、それぞれ異なる種類の情報であり、単純な性能ランキングへまとめられない。未確定の点は、次の公開更新や再現可能な評価で確かめるべき観測点として残る。
\end{claimboundary}
